\subsection{Change of Base Field}

We keep the notation of the previous subsection. Let \(K'\) be an algebraically closed field such that the element \(n \cdot 1\) of \(K'\) is not zero. The groups \(\mu_{n}(\mathrm{K})\) and \(\mu_{n}(\mathrm{K}')\) are cyclic of order n (V, §11, No.~2, p.~78, Theorem 1). Choose an isomorphism \(\varphi\) from \(\mu_{n}(\mathrm{K})\) to \(\mu_{n}(\mathrm{K}')\) . Let \(\pi\) be a linear representation of G in a finite-dimensional K-vector space, and let \(\pi'\) be a linear representation of G in a finite-dimensional \(K'\) -vector space. We say that \(\pi\) and \(\pi'\) are related (through \(\varphi\) ) if for every \(g \in G\) and \(\omega \in \mu_{n}(\mathrm{K})\) , the multiplicity of \(\omega\) as an eigenvalue of \(\pi(g)\) is equal to the multiplicity of \(\varphi(\omega)\) as an eigenvalue of \(\pi'(g)\) . When this is the case, \(\pi\) and \(\pi'\) have the same dimension, as can be seen by taking g = 1.

Let \(\pi_1\) and \(\pi_2\) (resp. \(\pi_1'\) and \(\pi_2'\) ) be linear representations of \(G\) in finite-dimensional vector spaces over \(K\) (resp. \(K'\) ). We have the following properties:

\begin{enumerate}
\def\labelenumi{\alph{enumi})}
\item
  If \(\pi_1\) is related to \(\pi_1'\) and \(\pi_2'\) , then \(\pi_1'\) and \(\pi_2'\) are isomorphic.
\item
  If \(\pi_1\) is related to \(\pi_1'\) and \(\pi_2\) to \(\pi_2'\) , then \(\pi_1 \oplus \pi_2\) is related to \(\pi_1' \oplus \pi_2'\) and \(\pi_1 \otimes \pi_2\) is related to \(\pi_1' \otimes \pi_2'\) .
\end{enumerate}

Assertion a) follows from Corollary 5 of VIII, p.~402, and assertion b) is clear.

Here, we denote by \(\mathcal{S}_{\mathrm{K}}(\mathrm{G})\) the set of classes of simple K{[}G{]}-modules, previously denoted by \(\widehat{G}\) , and we define \(\mathcal{S}_{\mathrm{K}^{\prime}}(\mathrm{G})\) likewise. The sets \(\mathcal{S}_{\mathrm{K}}(\mathrm{G})\) and \(\mathcal{S}_{\mathrm{K}^{\prime}}(\mathrm{G})\) are both finite, with cardinal the number of conjugacy classes (VIII, p.~411, Proposition 5).

PROPOSITION 10. --- There exists a unique mapping \(\varphi_{\mathrm{G}}\) from \(\mathcal{S}_{\mathrm{K}}(\mathrm{G})\) to \(\mathcal{S}_{\mathrm{K}'}(\mathrm{G})\) such that \(\lambda\) and \(\varphi_{\mathrm{G}}(\lambda)\) are related through \(\varphi\) for every \(\lambda\) in \(\mathcal{S}_{\mathrm{K}}(\mathrm{G})\) . Moreover, \(\varphi_{\mathrm{G}}\) is bijective.

The uniqueness of \(\varphi_{G}\) follows from property a) above.

\textbf{A) Suppose that the field K has characteristic 0.}

The group \(\mu_n(\mathrm{K})\) is cyclic (V, §11, No.~2, p.~78, Theorem 1); choose a generator \(\zeta\) of this group. Consider the ring homomorphism \(\rho: \mathbf{Z}[\mathrm{X}] \to \mathcal{O}_n\) that sends X to \(\zeta\) . It is surjective. The cyclotomic polynomial \(\Phi_n(\mathrm{X})\) is irreducible in \(\mathbf{Q}[\mathrm{X}]\) (V, §11, No.~5, p.~84, Theorem 2); it is therefore the minimal polynomial of \(\zeta\) over \(\mathbf{Q}\) . The polynomial \(\Phi_n\) is monic with integer coefficients (V, §11, No.~4, p.~81). Let \(\mathrm{P} \in \mathbf{Z}[\mathrm{X}]\) be a polynomial such that \(\mathrm{P}(\zeta) = 0\) ; by Euclidean division of polynomials (IV, §1, No.~6, p.~10), there exist two polynomial Q and R in \(\mathbf{Z}[\mathrm{X}]\) such that \(\mathrm{P} = \mathrm{Q}\Phi_n + \mathrm{R}\) and \(\deg(\mathrm{R}) < \deg(\Phi_n)\) . We have \(\mathrm{R}(\zeta) = 0\) , and therefore \(\mathrm{R} = 0\) because \(\Phi_n\) is the minimal polynomial of \(\zeta\) . Consequently, the kernel of \(\rho\) is the ideal \(\Phi_n\mathbf{Z}[\mathrm{X}]\) of \(\mathbf{Z}[\mathrm{X}]\) , and \(\rho\) induces a ring isomorphism from \(\mathbf{Z}[\mathrm{X}] / \Phi_n\mathbf{Z}[\mathrm{X}]\) to \(\mathcal{O}_n\) .

Set \(\zeta' = \varphi(\zeta)\) ; it is a primitive \(n\) -th root of unity in \(\mathbf{K}'\) , and we therefore have \(\Phi_n(\zeta') = 0\) (V, §11, No.~5, p.~83, Lemma 3). Consequently, there exists a homomorphism \(\varphi_0\) from the ring \(\mathcal{O}_n\) to the field \(\mathbf{K}'\) that transforms \(\zeta\) into \(\zeta'\) ; it extends the mapping \(\varphi\) from \(\mu_n(\mathrm{K})\) to \(\mu_n(\mathrm{K}')\) . Let \(\mathcal{O}\) be the subring of \(\mathrm{K}\) consisting of the elements \(\frac{a}{n^r}\) with \(a \in \mathcal{O}_n\) and \(r \in \mathbf{N}\) . Since \(n \cdot 1\) is invertible in \(\mathbf{K}'\) , the homomorphism \(\varphi_0\) extends to a homomorphism \(\varphi_1\) from \(\mathcal{O}\) to \(\mathbf{K}'\) .

We identify the algebra \(\mathcal{O}[\mathrm{G}]\) of the group \(\mathrm{G}\) over \(\mathcal{O}\) with a subring of the algebra \(\mathrm{K}[\mathrm{G}]\) and define a ring homomorphism \(\Phi\) from \(\mathcal{O}[\mathrm{G}]\) to \(\mathrm{K}'[\mathrm{G}]\) by the formula

\[
\Phi \bigg (\sum_ {g \in \mathrm{G}} a _ {g} g \bigg) = \sum_ {g \in \mathrm{G}} \varphi_ {1} (a _ {g}) g.\tag{42}
\]

Denote by \(\mathcal{C}\) the set of conjugacy classes of G. For C in \(\mathcal{C}\) , denote by \(u_{\mathrm{C}}\) the element \(\sum_{g\in \mathrm{C}}g\) of \(\mathcal{O}[\mathrm{G}]\) ; the family \((u_{\mathrm{C}})_{\mathrm{C}\in \mathcal{C}}\) is a basis over \(\mathcal{O}\) of the center \(\mathrm{Z}(\mathcal{O}[\mathrm{G}])\) of the algebra \(\mathcal{O}[\mathrm{G}]\) (VIII, p.~398). For any element \(\lambda\) of \(\mathcal{S}_{\mathrm{K}}(\mathrm{G})\) , denote by \(\chi_{\lambda}\) its character, by \(d_{\lambda}\) its dimension, and by \(e_{\lambda}\) the

element of K{[}G{]} defined by

\[
e _ {\lambda} = | \mathrm{G} | ^ {- 1} d _ {\lambda} \sum_ {g \in \mathrm{G}} \chi_ {\lambda} (g ^ {- 1}) g.\tag{43}
\]

The family \((e_{\lambda})\) is a basis over K of the center \(Z(K[G])\) of the ring K{[}G{]} (VIII, p.~408). We have

\[
e _ {\lambda} = \sum_ {\mathrm{C} \in \mathcal {C}} \alpha_ {\lambda , \mathrm{C}} u _ {\mathrm{C}}\tag{44}
\]

with

\[
\alpha_ {\lambda , \mathrm{C}} = | \mathrm{G} | ^ {- 1} d _ {\lambda} \chi_ {\lambda} (\mathrm{C} ^ {- 1}).\tag{45}
\]

For \(C \in \mathcal{C}\) and \(\lambda \in \mathcal{S}_{\mathrm{K}}(\mathrm{G})\) , set \(\beta_{\mathrm{C},\lambda} = |\mathrm{G}| d_{\lambda}^{-1} d(\mathrm{C})^{-1} \chi_{\lambda}(\mathrm{C})\) . The entries of the matrix \((\alpha_{\lambda,\mathrm{C}})\) are in \(\mathcal{O}\), and it follows from formula (31) of VIII, p.~411 that its inverse matrix is the matrix \((\beta_{\mathrm{C},\lambda})\) , whose entries are also in \(\mathcal{O}\) by Proposition 9 of VIII, p.~415. Consequently, the family \((e_{\lambda})\) is a basis of the \(\mathcal{O}\)-module \(\mathrm{Z}(\mathcal{O}[\mathrm{G}])\) .

The elements \(\Phi(u_{\mathrm{C}}) = \sum_{g\in \mathrm{C}}g\) of \(K^{\prime}[G]\) form a basis over \(K^{\prime}\) of the center \(Z(K^{\prime}[G])\) of the ring \(K^{\prime}[G]\) . We have

\[
\Phi (e _ {\lambda}) = \sum_ {\mathrm{C} \in \mathcal {C}} \varphi_ {1} (\alpha_ {\lambda , \mathrm{C}}) \Phi (u _ {\mathrm{C}}),\tag{46}
\]

and the matrix with entries \(\varphi_{1}(\alpha_{\lambda,\mathrm{C}})\) is invertible. The family of the \(\Phi(e_{\lambda})\) is therefore a basis of \(\mathrm{Z}(\mathrm{K}^{\prime}[\mathrm{G}])\) . The family \((e_{\lambda})\) is a partition of the idempotent 1 in \(\mathrm{Z}(\mathrm{K}[\mathrm{G}])\) (VIII, p.~146 and p.~408); in other words, we have

\[
\sum_ {\lambda} e _ {\lambda} = 1, \quad e _ {\lambda} ^ {2} = e _ {\lambda}, \quad e _ {\lambda} e _ {\mu} = 0 \quad \mathrm{if} \quad \lambda \neq \mu .
\]

It follows that the family of the \(\Phi(e_{\lambda})\) is a partition of the idempotent 1 in \(Z(K^{\prime}[G])\) ; since this family is a basis of \(Z(K^{\prime}[G])\) over \(K^{\prime}\) , its elements are the indecomposable idempotents in \(Z(K^{\prime}[G])\) (VIII, p.~148, Remark 4).

For \(\lambda'\) in \(\mathcal{S}_{\mathrm{K}'}(\mathrm{G})\) , define \(\chi_{\lambda'}, d_{\lambda'}\) , and \(e_{\lambda'}\) as above. By VIII, p.~408, the elements \(e_{\lambda'}\) are the indecomposable idempotents in \(\mathrm{Z}(\mathrm{K}'[\mathrm{G}])\) . Hence there exists a bijection \(\varphi_{\mathrm{G}}\) from \(\mathcal{S}_{\mathrm{K}}(\mathrm{G})\) to \(\mathcal{S}_{\mathrm{K}'}(\mathrm{G})\) such that \(\Phi(e_{\lambda}) = e_{\varphi_{\mathrm{G}}(\lambda)}\) for every \(\lambda\) in \(\mathcal{S}_{\mathrm{K}}(\mathrm{G})\) .

Let \(\lambda\) in \(\mathcal{S}_{\mathrm{K}}(\mathrm{G})\) ; set \(\lambda' = \varphi_{\mathrm{G}}(\lambda)\) . Let \((\mathrm{V}_{\lambda}, \pi_{\lambda})\) (resp. \((\mathrm{V}_{\lambda'}, \pi_{\lambda'})\) ) be a linear representation of G whose associated K{[}G{]}-module (resp. K'{[}G{]}-module) has class \(\lambda\) (resp. \(\lambda'\) ). Let us prove that \(\lambda\) and \(\lambda'\) are related. Let \(g\) be an element of G. Let \(\delta(\mathrm{T})\) be the determinant of the endomorphism \(1 + \mathrm{T}\pi_{\lambda}(g)\)

of the K{[}T{]}-module K{[}T{]} \(\otimes_{\mathrm{K}}\mathrm{V}_{\lambda}\) . Let \(\omega_{1},\ldots ,\omega_{d_{\lambda}}\) be the eigenvalues of \(\pi_{\lambda}(g)\) ; we have

\[
\delta (\mathrm{T}) = (1 + \mathrm{T} \omega_ {1}) \dots (1 + \mathrm{T} \omega_ {d _ {\lambda}}).
\]

We define \(\delta'(T)\) likewise, and we denote the eigenvalues of \(\pi_{\lambda'}(g)\) by \(\omega_1', \ldots, \omega_{d_{\lambda'}}'\) . The \(\mathcal{O}\) -module \(\mathcal{O}[G]\) is free with basis \(G\) , and the \(K\) -vector space \(K[G]\) has basis \(G\) . Denote by \(\Delta(T)\) the determinant of multiplication by \(1 + e_{\lambda}gT\) in the \(\mathcal{O}[T]\) -module \(\mathcal{O}[T] \otimes_{\mathcal{O}} \mathcal{O}[G]\) . It is also the determinant of multiplication by \(1 + e_{\lambda}gT\) in the \(K[T]\) -module \(K[T] \otimes_K K[G]\) . Let \(\overline{\varphi}_1\) be the homomorphism from \(\mathcal{O}[T]\) to \(K'[T]\) that extends \(\varphi_1\) and sends \(T\) to \(T\) . Since \(G\) is a basis of the \(K'\) -vector space \(K'[G]\) , the polynomial \(\overline{\varphi}_1(\Delta(T))\) is equal to the determinant \(\Delta'(T)\) of multiplication by \(1 + e_{\lambda'}gT\) in the \(K'\) -vector space \(K'[G]\) .

The algebra K{[}G{]} is the direct sum of its simple components \(e_{\mu}K[G]\) for \(\mu\) running through \(\mathcal{S}_{\mathrm{K}}(\mathrm{G})\) . For \(\mu\) different from \(\lambda\) , the element \(e_{\lambda}g\) annihilates \(e_{\mu}K[G]\) . Moreover, multiplication by \(e_{\lambda}g\) coincides with multiplication by g in \(e_{\lambda}K[G]\) . In view of VIII, p.~409 and Example 6 of VIII, p.~400, the representation of G in \(e_{\lambda}K[G]\) is the direct sum of \(d_{\lambda}\) representations of class \(\lambda\) . We consequently have \(\Delta(\mathrm{T}) = \delta(\mathrm{T})^{d_{\lambda}}\) .

Analogously, we have \(\Delta^{\prime}(\mathrm{T}) = \delta^{\prime}(\mathrm{T})^{d_{\lambda^{\prime}}}\) .

From the relation \(\Delta'(\mathrm{T}) = \overline{\varphi}_1(\Delta(\mathrm{T}))\) , we deduce first that \(d_{\lambda}^2 = d_{\lambda'}^2\) , and therefore \(d_{\lambda} = d_{\lambda'}\) , and then that the sequence \(\varphi(\omega_1), \ldots, \varphi(\omega_{d_\lambda})\) can be deduced from the sequence \((\omega_1', \ldots, \omega_{d_{\lambda'}'}')\) by a permutation of the set of indices.

Since this is true for every element \(g\) of \(G\) , the representations \(\lambda\) and \(\lambda'\) are related. We have therefore proved Proposition 10 when the field \(K\) has characteristic 0.

\textbf{B) General case.}

Let L be an algebraically closed field of characteristic 0 (for example, an algebraic closure of Q). Denote by \(\mathcal{S}_{\mathrm{L}}(\mathrm{G})\) the set of classes of simple L{[}G{]}-modules. Choose an isomorphism \(\eta\) from the group \(\mu_{n}(\mathrm{L})\) to the group \(\mu_{n}(\mathrm{K})\) , and set \(\eta' = \varphi \circ \eta\) . By part A) of the proof, there exist bijections

\[
\eta_ {\mathrm{G}}: \mathcal {S} _ {\mathrm{L}} (\mathrm{G}) \to \mathcal {S} _ {\mathrm{K}} (\mathrm{G}), \quad \eta_ {\mathrm{G}} ^ {\prime}: \mathcal {S} _ {\mathrm{L}} (\mathrm{G}) \to \mathcal {S} _ {\mathrm {K^ {\prime}}} (\mathrm{G})
\]

with the following property: for every \(\lambda\) in \(\mathcal{S}_{\mathrm{L}}(\mathrm{G})\) , the representations \(\lambda\) and \(\eta_{\mathrm{G}}(\lambda)\) are related through \(\eta\) , and the representations \(\lambda\) and \(\eta_{\mathrm{G}}^{\prime}(\lambda)\) are related through \(\eta^{\prime}\) . The bijection \(\varphi_{G} = \eta_{G}^{\prime} \circ \eta_{G}^{-1}\) has the desired properties.

The bijection \(\varphi_{G}\) from \(\mathcal{S}_{K}(G)\) to \(\mathcal{S}_{K'}(G)\) extends to an isomorphism, also denoted by \(\varphi_{G}\) , from the Grothendieck group \(R_{K}(G)\) to the group \(R_{K'}(G)\) .

Remark 1. --- Suppose that \(K'\) is an extension of K and that the isomorphism \(\varphi\) is the mapping \(\xi \mapsto \xi \cdot 1\) ; then the mapping \(\varphi_{G}\) is given by extension of scalars from K to \(K'\) .

COROLLARY 1. --- The mapping \(\varphi_{\mathrm{G}}\) is a ring isomorphism from \(\mathrm{R}_{\mathrm{K}}(\mathrm{G})\) to \(\mathrm{R}_{\mathrm{K}^{\prime}}(\mathrm{G})\) . For every finite-dimensional representation \(\pi\) of \(\mathrm{G}\) in a K-vector space, we have \(\varphi_{\mathrm{G}}([\pi]) = [\pi^{\prime}]\) , where \(\pi^{\prime}\) is a representation related to \(\pi\) through \(\varphi\) .

This follows from the semisimplicity of the representations of G and property b) of VIII, p.~416.

COROLLARY 2. --- The dimension of each simple representation of G divides the order of G.

This follows from Proposition 10 and Proposition 9 of VIII, p.~415.

Remarks. --- 2) Suppose that the group G is abelian. We saw in the remark of VIII, p.~414 that \(\mathcal{S}_{\mathrm{K}}(\mathrm{G})\) can be identified with the set \(\operatorname{Hom}(\mathrm{G}, \mu_n(\mathrm{K}))\) . Likewise, \(\mathcal{S}_{\mathrm{K}'}(\mathrm{G})\) can be identified with \(\operatorname{Hom}(\mathrm{G}, \mu_n(\mathrm{K}'))\) . With these identifications, the bijection \(\varphi_{\mathrm{G}}\) is simply the mapping \(\chi \mapsto \varphi \circ \chi\) .

\begin{enumerate}
\def\labelenumi{\arabic{enumi})}
\setcounter{enumi}{2}
\tightlist
\item
  Let \(\pi_1\) and \(\pi_2\) be linear representations of G in finite-dimensional vector spaces over K. For \(i = 1,2\) , let \(\pi_i'\) be a representation related to \(\pi_i\) through \(\varphi\) . We have
\end{enumerate}

\[
\dim_ {K} \operatorname{Hom} _ {K} (\pi_ {1}, \pi_ {2}) = \dim_ {K ^ {\prime}} \operatorname{Hom} _ {K ^ {\prime}} (\pi_ {1} ^ {\prime}, \pi_ {2} ^ {\prime}).\tag{47}
\]

The proof follows that of Corollary VIII, p.~410, by reducing to the case when the \(\pi_{i}\) (and therefore the \(\pi_{i}^{\prime}\) ) are simple.

\begin{enumerate}
\def\labelenumi{\arabic{enumi})}
\setcounter{enumi}{3}
\tightlist
\item
  Let H be a subgroup of G of cardinal m. The isomorphism \(\varphi\) restricts to an isomorphism from \(\mu_{m}(\mathrm{K})\) to \(\mu_{m}(\mathrm{K}^{\prime})\) and, consequently, a ring isomorphism \(\varphi_{H}\) from \(\mathrm{R}_{\mathrm{K}}(\mathrm{H})\) to \(\mathrm{R}_{\mathrm{K}^{\prime}}(\mathrm{H})\) . The following diagrams commute:
\end{enumerate}

\[
\begin{array}{c c} \mathrm {R_ {K} (G)} \xrightarrow {\mathrm {Res_ {H} ^ {G}}} \mathrm {R_ {K} (H)} & \qquad \qquad \mathrm {R_ {K} (H)} \xrightarrow {\mathrm {Ind_ {H} ^ {G}}} \mathrm {R_ {K} (G)} \\ \Biggl \downarrow \varphi_ {\mathrm{G}} & \Biggl \downarrow \varphi_ {\mathrm{H}} \\ \mathrm {R_ {K^ {\prime}} (G)} \xrightarrow {\mathrm {Res_ {H} ^ {G}}} \mathrm {R_ {K^ {\prime}} (H)}, & \qquad \qquad \mathrm {R_ {K^ {\prime}} (H)} \xrightarrow {\mathrm {Ind_ {H} ^ {G}}} \mathrm {R_ {K^ {\prime}} (G)}. \end{array}
\]

The commutativity of the first diagram is obvious, and that of the second follows from it using Frobenius reciprocity and formula (47).

\begin{enumerate}
\def\labelenumi{\arabic{enumi})}
\setcounter{enumi}{4}
\tightlist
\item
  Suppose that \(G\) is the product \(G' \times G''\) of two finite groups. We define isomorphisms \(\varphi_{G'}\) and \(\varphi_{G''}\) as in the previous example.
\end{enumerate}

We have a commutative diagram

\[
\begin{array}{c} \mathrm{R} _ {\mathrm{K}} (\mathrm{G} ^ {\prime}) \otimes_ {\mathbf {Z}} \mathrm{R} _ {\mathrm{K}} (\mathrm{G} ^ {\prime \prime}) \xrightarrow {\kappa} \mathrm{R} _ {\mathrm{K}} (\mathrm{G}) \\ \Biggl \downarrow \varphi_ {\mathrm{G} ^ {\prime}} \otimes \varphi_ {\mathrm{G} ^ {\prime \prime}} \qquad \qquad \qquad \Biggl \downarrow \varphi_ {\mathrm{G}} \\ \mathrm{R} _ {\mathrm{K} ^ {\prime}} (\mathrm{G} ^ {\prime}) \otimes_ {\mathbf {Z}} \mathrm{R} _ {\mathrm{K} ^ {\prime}} (\mathrm{G} ^ {\prime \prime}) \xrightarrow {\kappa^ {\prime}} \mathrm{R} _ {\mathrm{K} ^ {\prime}} (\mathrm{G})  , \end{array}
\]

where the isomorphisms \(\kappa\) and \(\kappa'\) are those defined in VIII, p.~406.

